\documentclass{article}
\iffalse
use pdflatex to compile and xdg-open to open pdf, make images as cbr files
mathbb{R} is set
\fi
\usepackage{amsmath}
\usepackage{amsfonts}
\usepackage{xcolor}
% following inclusions should allow graphs
\usepackage{tikz}
\usetikzlibrary{graphs, graphdrawing}
\usetikzlibrary{positioning}
\usegdlibrary{force}
\begin{document}
This document is a showcase for using a graphing library\\
compile as `lualatex graphs_showcase.tex`\\
1. using \\graph:\\
\begin{tikzpicture}
	\graph [spring layout, nodes={draw, circle}, node distance=1.5cm] {
		A -- {B, C, D},
		B -- {C, E},
		C -- D,
		D -- E
	};
\end{tikzpicture}

2. specifying circle-node positions :\\
\begin{tikzpicture}
    \node[circle,draw] (A) at (0,0) {A};
    \node[circle,draw] (B) at (2,1) {B};
    \node[circle,draw] (C) at (2,-1) {C};
    \node[circle,draw] (D) at (4,0) {D};

    \draw (A) -- (B);
    \draw (A) -- (C);
    \draw (B) -- (D);
    \draw (C) -- (D);
\end{tikzpicture}

3. Hesse diagram with specified positions verbally and nodes are sets:\\
\begin{tikzpicture}
    \node (empty) at (0,0) {$\emptyset$};
    
    \node (1)  [above left=of empty]  {$\{1\}$};
    \node (2)  [above=of empty]  {$\{2\}$};
    \node (3)  [above right=of empty]  {$\{3\}$};
    
    \node (12) [above=of 1]  {$\{1,2\}$};
    \node (13) [above=of 2]  {$\{1,3\}$};
    \node (23) [above=of 3]  {$\{2,3\}$};
    
    \node (123) [above=of 13] {$\{1,2,3\}$};
    
    \draw (empty) -- (1);
    \draw (empty) -- (2);
    \draw (empty) -- (3);
    
    \draw (1) -- (12);
    \draw (1) -- (13);
    \draw (2) -- (12);
    \draw (2) -- (23);
    \draw (3) -- (13);
    \draw (3) -- (23);
    
    \draw (12) -- (123);
    \draw (13) -- (123);
    \draw (23) -- (123);
\end{tikzpicture}

\end{document}
